\begin{enumerate}
    \item \textbf{Delimitación del perfil de participantes:} Se recomienda replicar el diseño experimental incorporando grupos compuestos por jugadores frecuentes (\textit{gamers}) y profesionales de la industria del desarrollo de videojuegos. Dado que la muestra del presente estudio estuvo integrada por estudiantes universitarios, la extensión a públicos especializados e inmersos en el sector productivo permitirá determinar la intensidad y la persistencia del sesgo de etiquetado en actores directamente vinculados con el ecosistema de los videojuegos.
    \item \textbf{Implementación de estudios longitudinales de percepción:} Es aconsejable diseñar protocolos de evaluación con exposiciones prolongadas al videojuego, en lugar de sesiones únicas de prueba. Evaluar la percepción de los usuarios a lo largo del tiempo permitirá analizar si la brecha de valoración causada por la etiqueta de autoría se mantiene constante, se amplifica o se neutraliza a medida que el jugador se enfoca de forma sostenida en la experiencia lúdica y en la maestría de las mecánicas.
    \item \textbf{Complementación con métricas objetivas y biométricas:} Para reducir la subjetividad inherente a los cuestionarios autoinformados como el PXI, se sugiere integrar herramientas de medición fisiológica o de comportamiento en tiempo real, tales como seguimiento ocular (\textit{eye-tracking}), análisis de frecuencia cardíaca o registro detallado de interacciones por minuto (\textit{input logging}). Esto ofrecerá una capa de validación cruzada que aísle con mayor precisión la respuesta cognitiva del jugador frente a la calidad técnica del código.
    \item \textbf{Escalamiento a arquitecturas y entornos de mayor complejidad:} Se recomienda aplicar el protocolo de pruebas a proyectos de software interactivo que superen el alcance de prototipos 2D, integrando entornos tridimensionales, arquitecturas de código distribuido y sistemas de física compleja en motores como Godot o similares. Evaluar la generación asistida de código en escenarios de alta densidad algorítmica permitirá verificar si la paridad técnica observada frente al desarrollo humano se sostiene cuando aumentan la exigencia computacional y la interdependencia entre módulos.
    \item \textbf{Estandarización de flujos de trabajo híbrido y pautas de integración:} Para los equipos de desarrollo y la industria de videojuegos, se sugiere formalizar guías de buenas prácticas para la adopción de herramientas de IA generativa en la etapa de programación. Dichas pautas deben centrarse en el uso de la IA como un asistente de refactorización, optimización y generación de prototipos lógicos, asegurando que la revisión de código por parte de programadores humanos garantice los estándares de seguridad, mantenibilidad y rendimiento del proyecto.
    \item \textbf{Estrategias de comunicación y divulgación de autoría:} Ante la evidencia del sesgo de percepción asociado al etiquetado de la IA, se aconseja que las empresas y desarrolladores diseñen estrategias transparentes pero contextualizadas al informar la incorporación de tecnologías generativas en sus productos. La comunicación debe enfatizar la calidad técnica, el rigor de la validación humana y la experiencia funcional del usuario por encima del mero etiquetado de las herramientas utilizadas, evitando que prejuicios o estigmas infundados opaquen el valor real del software.
    \item \textbf{Ampliación del diseño metodológico y condiciones de prueba:} Se propone que futuras investigaciones incrementen el tamaño muestral e incorporen una nueva condición experimental basada en la información parcial de la autoría. En esta modalidad, los participantes sabrán que uno de los prototipos fue desarrollado con asistencia de IA, pero desconocerán a cuál de ellos corresponde. Este enfoque permitirá evaluar la capacidad real de los usuarios para discriminar código e interfaces generadas de forma automatizada frente a la mera manifestación de un sesgo inducido por la etiqueta explícita.
\end{enumerate}